\documentclass[12pt,a4paper]{article}

\usepackage[margin=2cm]{geometry}
\usepackage[T1]{fontenc}
\usepackage{lmodern}
\usepackage{xcolor}

\setlength{\parindent}{0pt}
\setlength{\parskip}{0.7em}

\definecolor{TitleBlue}{HTML}{0D617C}

\title{Ultra-Short English Memorization Script\\[4pt]
\large Cross-Dataset Retinal Vessel Segmentation with Lightweight Retinal-Specific Priors}
\author{COMP5423 Final Project}
\date{April 21, 2026}

\begin{document}

\maketitle

This document is a short English-only speaking script for the 12-slide presentation.

Each slide keeps only 3 to 5 short sentences, so it is easier to memorize quickly.

\newpage

\section*{Slide 1: Title and Main Claim}

Good morning everyone.

This paper studies retinal vessel segmentation, which means detecting blood vessels in images of the eye.

Cross-dataset means we train on one dataset and test on another.

The main message is simple: better preprocessing helps more than a larger model.

\newpage

\section*{Slide 2: Why Same-Dataset Scores Are Not Enough}

This paper matters because medical images from different cameras do not look the same.

This difference is called domain shift.

Thin vessels are hard to detect because they are small and low in contrast.

So a model that works on one dataset may fail on another.

\newpage

\section*{Slide 3: Strict Bidirectional Setup}

The paper uses two public datasets: STARE and HRF.

It tests both directions: train on STARE and test on HRF, and train on HRF and test on STARE.

This is called bidirectional transfer.

The setup is strict because there is no target-domain fine-tuning.

\newpage

\section*{Slide 4: Method and Four Recipes}

The backbone is a standard U-Net, which is a common segmentation network.

The paper compares simple recipes instead of changing the network.

The key preprocessing uses the green channel and histogram equalization to improve contrast.

The full version also adds balanced loss, stronger augmentation, AdamW, and threshold search.

\newpage

\section*{Slide 5: Main Results}

The main result is clear.

From STARE to HRF, Green+Equalize raises Dice from 0.2109 to 0.5789.

From HRF to STARE, Full improved is best at 0.5430, but only slightly above Green+Equalize.

So preprocessing creates most of the gain.

\newpage

\section*{Slide 6: What Each Component Does}

Now we ask what each component really does.

Preprocessing gives the strongest and most consistent improvement.

Balanced loss helps with class imbalance, but it is weaker on the hardest cases.

Full improved is useful, but its extra gain is smaller than the jump from preprocessing.

\newpage

\section*{Slide 7: Training Dynamics and Transport Gap}

The training curves show that fast learning is not the same as better transfer.

In one direction, Full improved rises faster at first, but Green+Equalize finishes higher.

The paper also measures transport gap, which is the difference between source validation and target test performance.

A smaller gap means model selection is more stable, and Full improved is better on that point.

\newpage

\section*{Slide 8: Robustness Beyond the Average}

Average scores do not tell the full story.

The paper ranks images from hard to easy and looks at the bottom quartile, which means the hardest 25 percent.

Balanced only looks good in the middle, but it still fails badly on some hard images.

Preprocessing-heavy methods recover hard cases more reliably.

\newpage

\section*{Slide 9: Qualitative Case One}

This page shows qualitative results, which means visual examples instead of only numbers.

On image 12\_h, Green+Equalize restores many side branches that the baseline misses.

The mask also stays relatively clean.

This supports the idea that preprocessing already does most of the important work.

\newpage

\section*{Slide 10: Qualitative Case Two}

The image im0324 is a strong rescue case.

The baseline almost fails, but preprocessing brings back much more of the vessel tree.

The image im0077 is a counterexample, which means a case that shows a limit.

There, Balanced only looks cleaner, while stronger methods can over-segment.

\newpage

\section*{Slide 11: Limits and Boundaries}

The paper is useful, but it has clear limits.

It uses only two datasets, and not all training budgets are perfectly matched.

All images are resized to 256 by 256, which can hurt very thin vessels.

So the claim is careful: in this small setting, retinal priors matter more than a fancier model.

\newpage

\section*{Slide 12: Final Takeaways}

I will end with four takeaways.

First, retinal-specific preprocessing is the main reason performance improves.

Second, balanced loss helps, but it is not the main driver.

Third, real generalization needs more than one mean Dice score.

Before choosing a larger backbone, we should first align the input distribution.

\end{document}